% from ega2-1-fr.tex
\oldpage[II]{172}
qui est une $\mathcal{O}_X$-Algèbre graduée quasi-cohérente, la multiplication
y étant définie au moyen des homomorphismes canoniques
(\hyperref[II.3.2.6.1-fr]{3.2.6.1})
\[
  \mathcal{O}_X(m)\otimes_{\mathcal{O}_X}\mathcal{O}_X(n)
  \to\mathcal{O}_X(m+n)
\]
dont l'associativité est garantie par le diagramme commutatif
(\hyperref[II.2.5.11.4-fr]{2.5.11.4}). Prenons pour $\mathcal{S}'$ la
sous-$\mathcal{O}_X$-Algèbre graduée quasi-cohérente
$\mathcal{S}_X^{\geq}=\bigoplus_{n\geq0}\mathcal{O}_X(n)$ de
$\mathcal{S}_X$, à degrés positifs.

Enfin, nous considérerons le $q$-morphisme canonique
\begin{equation}
\label{II.8.6.1.2-fr}
  \alpha:\mathcal{S}\to\mathcal{S}_X^{\geq}
\tag{8.6.1.2}
\end{equation}
défini dans (\hyperref[II.3.3.2.3-fr]{3.3.2.3}) comme un homomorphisme
$\mathcal{S}\to q_*(\mathcal{S}_X)$, mais qui applique évidemment
$\mathcal{S}$ dans $q_*(\mathcal{S}_X^{\geq})$. Nous poserons
\begin{equation}
\label{II.8.6.1.3-fr}
  C_X=\operatorname{Spec}(\mathcal{S}_X^{\geq}),\qquad
  \widehat{C}_X=\operatorname{Proj}(\mathcal{S}_X^{\geq}[z]),\qquad
  X'=\operatorname{Proj}(\mathcal{S}_X^{\geq})
\tag{8.6.1.3}
\end{equation}
et nous désignerons par $E_X$ et $\widehat{E}_X$ le cône épointé affine et le
cône épointé projectif correspondants~; on notera
$\varepsilon_X:X\to C_X$, $i_X:C_X\to\widehat{C}_X$,
$j_X:X'\to\widehat{C}_X$, $p_X:\widehat{E}_X\to X'$,
$\pi_X:E_X\to X'$ les morphismes canoniques définis dans
(\hyperref[subsection:II.8.3-fr]{8.3}).
\end{env}

\begin{proposition}[8.6.2]
\label{II.8.6.2-fr}
Le morphisme structural $u:X'\to X$ est un isomorphisme, et le morphisme
$\operatorname{Proj}(\alpha)$ est partout défini et identique à $u$. Le
morphisme
$\operatorname{Proj}(\widehat{\alpha}):\widehat{C}_X\to\widehat{C}$ est
partout défini, et ses restrictions à $\widehat{E}_X$ et $E_X$ sont des
isomorphismes sur $\widehat{E}$ et $E$ respectivement. Enfin, lorsqu'on
identifie $X'$ et $X$ au moyen de $u$, les morphismes $p_X$ et $\pi_X$
s'identifient aux morphismes structuraux des $X$-préschémas $\widehat{E}_X$ et
$E_X$.
\end{proposition}

\begin{proof}
On peut évidemment se borner au cas où $Y=\operatorname{Spec}(A)$ est affine
et $\mathcal{S}=\widetilde{S}$~; alors $X$ est réunion des ouverts affines
$X_f$, où $f$ parcourt l'ensemble des éléments homogènes de $S_+$, l'anneau
de $X_f$ étant $S_{(f)}$. Il résulte en outre de
(\hyperref[II.8.2.7.1-fr]{8.2.7.1}) que l'on a
\begin{equation}
\label{II.8.6.2.1-fr}
  \Gamma(X_f,\mathcal{S}_X^{\geq})=S_f^{\geq}.
\tag{8.6.2.1}
\end{equation}

On a donc $u^{-1}(X_f)=\operatorname{Proj}(S_f^{\geq})$. Mais si
$f\in S_d$ ($d>0$), $\operatorname{Proj}(S_f^{\geq})$ est canoniquement
isomorphe à $\operatorname{Proj}((S_f^{\geq})^{(d)})$
(\hyperref[II.2.4.7-fr]{2.4.7}), et d'autre part
$(S_f^{\geq})^{(d)}=(S^{(d)})_f^{\geq}$ s'identifie à $S_{(f)}[T]$
(\hyperref[II.2.2.1-fr]{2.2.1}) par l'application $T\to f/1$~; on en conclut
(\hyperref[II.3.1.7-fr]{3.1.7}) que le morphisme structural
$u^{-1}(X_f)\to X_f$ est un isomorphisme, d'où la première assertion. Pour
démontrer la seconde, remarquons que la restriction
$u^{-1}(X_f)\cap G(\alpha)\to X=\operatorname{Proj}(S)$ de
$\operatorname{Proj}(\alpha)$ correspond à l'application canonique
$x\to x/1$ de $S$ dans $S_f^{\geq}$
(\hyperref[II.2.6.2-fr]{2.6.2})~; on en déduit tout d'abord que
$G(\alpha)=X'$, puis, en tenant compte de ce que
$u^{-1}(X_f)=(u^{-1}(X_f))_{f/1}$, il résulte de
(\hyperref[II.2.8.1.1-fr]{2.8.1.1}) que l'image de $u^{-1}(X_f)$ par
$\operatorname{Proj}(\alpha)$ est contenue dans $X_f$, et la restriction de
$\operatorname{Proj}(\alpha)$ à $u^{-1}(X_f)$, considérée comme morphisme
dans $X_f=\operatorname{Spec}(S_{(f)})$, est bien identique à celle de $u$.
Enfin, la formule (\hyperref[II.8.3.5.4-fr]{8.3.5.4}), appliquée à $p_X$ au
lieu de $p$, montre que
$p_X^{-1}(u^{-1}(X_f))=
\operatorname{Spec}((S_f^{\geq})_{f/1}^{\leq})$, et cet ouvert est, d'après
(\hyperref[II.8.5.4.1-fr]{8.5.4.1}), l'image réciproque par
$\operatorname{Proj}(\widehat{\alpha})$ de
$p^{-1}(X_f)=\operatorname{Spec}(S_f^{\leq})$
(\hyperref[II.8.3.5.3-fr]{8.3.5.3}). Compte tenu de
(\hyperref[II.8.2.3.2-fr]{8.2.3.2}), la restriction de
$\operatorname{Proj}(\widehat{\alpha})$ à
$p_X^{-1}(u^{-1}(X_f))$ correspond à l'isomorphisme réciproque de
(\hyperref[II.8.2.7.2-fr]{8.2.7.2}), restreint à $S_f^{\leq}$, d'où le
troisième point~; la dernière assertion est évidente par définition.

